\documentclass[12pt]{extreport}
\usepackage{subfigure}
\usepackage[vietnamese]{babel}
\usepackage[utf8]{inputenc}
\usepackage[T5]{fontenc}
\usepackage[left=3.50cm, right=2.00cm, top=3.50cm, bottom=3.00cm]{geometry}
\usepackage{fancybox,graphicx}
\usepackage{mathrsfs}
\usepackage{amsfonts}
\usepackage{longtable,array}
\usepackage{multirow}
\newlength\mylength
\newcolumntype{C}[1]{>{\centering\arraybackslash}p{#1}}
\usepackage[intlimits]{amsmath}
\usepackage{array}
\usepackage[unicode]{hyperref}
\usepackage{algorithm}
\usepackage{algorithmicx}
\makeatletter
\renewcommand{\ALG@name}{Thuật toán}
\makeatother
\usepackage{algpseudocode}
\usepackage{amsxtra,amssymb,latexsym,amscd,amsthm}
\usepackage{enumitem}
\usepackage{tikz}
\usetikzlibrary{shapes.geometric}
\usetikzlibrary{positioning,automata}
\usepackage{scrextend}
\usepackage{setspace}
\usepackage{hyperref}
\usepackage{tabularx}
\usepackage{multirow}
\usepackage[a4paper, margin=1in]{geometry}
\usepackage{xltabular}
\usepackage{indentfirst}
\usepackage{float}

% Môi trường định lý
\newtheorem{pro}{Bài toán}
\newtheorem*{constr}{Ràng buộc}
\newtheorem*{calfunc}{Các hàm được thực thi}
\newtheorem*{Sol}{Giải thuật}
\newtheorem*{Anal}{Phân tích giải thuật}

% Header/Footer
\usepackage{fancyhdr}
\pagestyle{fancy}
\fancyhf{}
\rhead{Tìm hiểu và xây dựng ứng dụng trên kiến trúc Microservice}
\cfoot{\thepage}
\renewcommand{\headrulewidth}{0.4pt}
\renewcommand{\footrulewidth}{0pt}

% Code listing
\usepackage{xcolor}
\usepackage{listings}

\definecolor{mGreen}{rgb}{0,0.6,0}
\definecolor{mGray}{rgb}{0.5,0.5,0.5}
\definecolor{mPurple}{rgb}{0.58,0,0.82}
\definecolor{backgroundColour}{rgb}{0.95,0.95,0.92}
\definecolor{codeblue}{rgb}{0,0,0.8}
\definecolor{codegreen}{rgb}{0,0.6,0}
\definecolor{codegray}{rgb}{0.5,0.5,0.5}
\definecolor{codepurple}{rgb}{0.58,0,0.82}
\definecolor{backcolour}{rgb}{0.97,0.97,0.95}

\setlength{\parindent}{1.5em}

\lstdefinestyle{TypeScriptStyle}{
    backgroundcolor=\color{backcolour},
    commentstyle=\color{codegreen},
    keywordstyle=\color{magenta},
    numberstyle=\tiny\color{codegray},
    stringstyle=\color{codepurple},
    basicstyle=\footnotesize\ttfamily,
    breakatwhitespace=false,
    breaklines=true,
    captionpos=b,
    keepspaces=true,
    numbers=left,
    numbersep=5pt,
    showspaces=false,
    showstringspaces=false,
    showtabs=false,
    tabsize=2,
    language=JavaScript,
    morekeywords={interface,type,namespace,readonly,public,private,protected,async,await,implements,extends}
}

\lstdefinestyle{BashStyle}{
    backgroundcolor=\color{backcolour},
    commentstyle=\color{codegreen},
    keywordstyle=\color{blue},
    numberstyle=\tiny\color{codegray},
    stringstyle=\color{red},
    basicstyle=\footnotesize\ttfamily,
    breakatwhitespace=false,
    breaklines=true,
    captionpos=b,
    keepspaces=true,
    numbers=left,
    numbersep=5pt,
    showspaces=false,
    showstringspaces=false,
    showtabs=false,
    tabsize=2
}

\begin{document}

% ---- TRANG BÌA ----
\begin{titlepage}
    \centering
    \thispagestyle{empty}

    {\fontsize{13pt}{1}\selectfont\textbf{TRƯỜNG ĐẠI HỌC BÁCH KHOA HÀ NỘI}}\\
    {\fontsize{13pt}{1}\selectfont\textbf{VIỆN CÔNG NGHỆ THÔNG TIN VÀ TRUYỀN THÔNG}}\\
    \textbf{─────────── * ───────────}\\[1cm]

    \includegraphics[scale=0.5]{assets/logo_bk.png}\\[1.2cm]

    \textbf{\Large ĐỒ ÁN}\\[0.3cm]
    \textbf{\LARGE TỐT NGHIỆP ĐẠI HỌC}\\[0.3cm]
    \textbf{\large NGÀNH CÔNG NGHỆ THÔNG TIN}\\[1cm]

    \textbf{\Large XÂY DỰNG NỀN TẢNG DỊCH VỤ VI MÔ}\\[0.2cm]
    \textbf{\Large (MICROSERVICES) CHO GIAO DỊCH VÀ}\\[0.2cm]
    \textbf{\Large TRAO ĐỔI CÔNG VIỆC TRỰC TUYẾN}\\[2cm]

    \begin{flushleft}
        \hspace{2cm}\textbf{Sinh viên thực hiện: \hspace{0.3cm} Đinh Hoàng Hải Hảo}\\[0.2cm]
        \hspace{2cm}\textbf{Lớp\hspace{3.8cm}: \hspace{0.2cm} ATTT -- K59}\\[0.2cm]
        \hspace{2cm}\textbf{Giáo viên hướng dẫn: \hspace{0.2cm} [TS] Hoàng Văn Hiệp}\\
    \end{flushleft}

    \vfill

    {\large\textbf{HÀ NỘI 5-2026}}
\end{titlepage}

% ---- PHIẾU GIAO NHẬN ----
\chapter*{PHIẾU GIAO NHẬN ĐỒ ÁN TỐT NGHIỆP}
\addcontentsline{toc}{chapter}{PHIẾU GIAO NHẬN ĐỒ ÁN TỐT NGHIỆP}
\thispagestyle{empty}

\textbf{1. Thông tin về sinh viên}

Họ và tên sinh viên: Đinh Hoàng Hải Hảo

Điện thoại liên lạc: 0973408415 \hspace{1cm} Email: haodinh0123@gmail.com

Lớp: ATTT-K59 \hspace{2cm} Hệ đào tạo: Kỹ sư

Thời gian làm đồ án tốt nghiệp: Từ ngày 11/02/2026 đến 24/05/2026

\bigskip
\textbf{2. Mục đích nội dung của đồ án tốt nghiệp}
\begin{itemize}
    \item Tìm hiểu kiến trúc Microservice và các công nghệ cloud-native hiện đại
    \item Xây dựng nền tảng giao dịch dịch vụ/công việc (gigs) trực tuyến theo mô hình Microservice
    \item Triển khai hệ thống trên Kubernetes (K3s) với CI/CD tự động
\end{itemize}

\textbf{3. Các nhiệm vụ cụ thể của đồ án tốt nghiệp}
\begin{itemize}
    \item Thiết kế và xây dựng 8 microservices độc lập (Node.js/TypeScript)
    \item Xây dựng giao diện người dùng hiện đại với React/Vite
    \item Tích hợp hệ thống tìm kiếm Elasticsearch, nhắn tin real-time Socket.io, thanh toán Stripe
    \item Triển khai trên Kubernetes với Docker, CI/CD GitHub Actions
\end{itemize}

\textbf{4. Lời cam đoan của sinh viên:}

Tôi -- Đinh Hoàng Hải Hảo -- cam kết ĐATN là công trình nghiên cứu của bản thân tôi dưới sự hướng dẫn của TS. Hoàng Văn Hiệp.

Các kết quả nêu trong ĐATN là trung thực, không phải là sao chép toàn văn của bất kỳ công trình nào khác.

\begin{flushright}
    Hà Nội, ngày \ldots tháng \ldots năm \ldots\\[0.5cm]
    Tác giả ĐATN\\[1cm]
    \textit{Đinh Hoàng Hải Hảo}
\end{flushright}

\textbf{5. Xác nhận của giáo viên hướng dẫn} về mức độ hoàn thành của ĐATN và cho phép bảo vệ:

\begin{flushright}
    Hà Nội, ngày \ldots tháng \ldots năm \ldots\\[0.5cm]
    Giáo viên hướng dẫn\\[1cm]
    \textit{TS. Hoàng Văn Hiệp}
\end{flushright}

\cleardoublepage

% ---- LỜI CẢM ƠN ----
\chapter*{LỜI CẢM ƠN}
\addcontentsline{toc}{chapter}{LỜI CẢM ƠN}

Năm tháng trôi qua tựa như một cơn gió lướt qua tuổi thanh xuân. Năm năm gắn với Bách Khoa, không dài cũng không ngắn, nhưng đó sẽ mãi là ký ức của một thời tuổi trẻ khao khát và dại khờ.

Lời đầu tiên tôi xin gửi lời cảm ơn tới TS. Hoàng Văn Hiệp, thày giáo hướng dẫn đồ án này. Cảm ơn thày đã gợi ý đề tài cũng như đưa ra các giải pháp, hướng đi và đồng hành giúp tôi hoàn thành đồ án trong thời gian cho phép và đạt kết quả tốt nhất.

Bên cạnh đó tôi xin gửi lời cảm ơn đến gia đình và bạn bè đã luôn là chỗ dựa tinh thần và giúp đỡ trong thời gian làm đồ án cũng như suốt những năm tháng học tập và gắn bó với Bách Khoa. Cảm ơn tất cả thày cô giáo trong trường và đặc biệt là các thày cô trong viện đã đồng hành, dìu dắt em qua những năm tháng học tập tại ngôi trường Bách Khoa thân yêu.

Do thời gian làm đồ án có hạn và kiến thức cũng như trình độ còn hạn chế nên không thể tránh khỏi những thiếu sót. Vì vậy, sự đóng góp ý kiến, nhận xét của thày cô cũng như các bạn sinh viên là vô cùng quý giá giúp đồ án có thể hoàn thiện hơn!

\begin{flushright}
    \textit{Xin chân thành cảm ơn!}
\end{flushright}

\cleardoublepage

% ---- TÓM TẮT ----
\chapter*{TÓM TẮT NỘI DUNG ĐỒ ÁN TỐT NGHIỆP}
\addcontentsline{toc}{chapter}{TÓM TẮT NỘI DUNG ĐỒ ÁN TỐT NGHIỆP}

Trong bối cảnh công nghệ thông tin phát triển mạnh mẽ, các ứng dụng web hiện đại đòi hỏi khả năng mở rộng, độ sẵn sàng cao và tốc độ phản hồi nhanh. Kiến trúc nguyên khối (monolithic) truyền thống gặp nhiều hạn chế khi hệ thống lớn dần, dẫn đến khó khăn trong việc bảo trì, nâng cấp và triển khai.

Qua việc tìm kiếm giải pháp, tôi nhận thấy nhiều tập đoàn lớn như Amazon, Netflix, Uber đã giải quyết vấn đề này bằng kiến trúc microservices --- tập hợp các dịch vụ nhỏ độc lập, giao tiếp với nhau thông qua các giao thức nhẹ. Ý tưởng là chia nhỏ ứng dụng lớn thành các dịch vụ nhỏ kết nối với nhau, điều này làm cho ứng dụng trở nên dễ hiểu, dễ phát triển, thử nghiệm, triển khai và trở nên linh hoạt hơn.

Đồ án này thực hiện xây dựng \textbf{``ITHUST App''} --- một nền tảng dịch vụ vi mô (microservices) mở rộng, hỗ trợ giao dịch và trao đổi công việc/dịch vụ (gigs) với các tính năng từ xác thực, tìm kiếm thời gian thực, nhắn tin, thanh toán đến hệ thống đánh giá. Hệ thống được thiết kế cloud-native, hướng tới độ sẵn sàng cao và tốc độ phản hồi nhanh.

Hệ thống bao gồm \textbf{8 microservices} độc lập:
\begin{itemize}
    \item \textbf{Gateway Service}: Cổng API tập trung xử lý định tuyến, rate limiting và xác thực
    \item \textbf{Auth Service}: Quản lý đăng ký, đăng nhập, phân quyền và cấp phát JWT
    \item \textbf{Users Service}: Quản lý hồ sơ người mua/người bán
    \item \textbf{Gig Service}: Quản lý dịch vụ/công việc (gigs) với tìm kiếm Elasticsearch
    \item \textbf{Order Service}: Xử lý đặt hàng, tích hợp thanh toán Stripe
    \item \textbf{Chat Service}: Giao tiếp thời gian thực qua Socket.io
    \item \textbf{Review Service}: Đánh giá và phản hồi đơn hàng
    \item \textbf{Notification Service}: Gửi email và thông báo (broker-driven)
\end{itemize}

Công nghệ sử dụng bao gồm \textbf{Node.js} và \textbf{TypeScript} cho backend, \textbf{React/Vite} cho frontend, \textbf{MongoDB/MySQL/PostgreSQL} cho cơ sở dữ liệu, \textbf{Redis} cho caching, \textbf{RabbitMQ} cho message broker, \textbf{Elasticsearch} cho tìm kiếm, và \textbf{Kubernetes (K3s)} cho triển khai.

\cleardoublepage

% ---- ABSTRACT ----
\chapter*{ABSTRACT}
\addcontentsline{toc}{chapter}{ABSTRACT}

In the context of rapid information technology development, modern web applications require scalability, high availability, and fast response times. Traditional monolithic architecture faces many limitations as systems grow larger, leading to difficulties in maintenance, upgrades, and deployment.

Through researching solutions, I found that many large corporations such as Amazon, Netflix, and Uber have solved this problem using microservices architecture --- a collection of small, independent services that communicate with each other through lightweight protocols. The idea is to break large applications into smaller, interconnected services, making the application easier to understand, develop, test, deploy, and more flexible.

This thesis implements \textbf{``ITHUST App''} --- an extended microservices platform supporting transactions and exchange of jobs/services (gigs) with features ranging from authentication, real-time search, messaging, payment to rating systems. The system is designed cloud-native, aiming for high availability and fast response speed.

The system includes \textbf{8 independent microservices}:
\begin{itemize}
    \item \textbf{Gateway Service}: Centralized API gateway handling routing, rate limiting, and authentication
    \item \textbf{Auth Service}: Managing registration, login, authorization, and JWT issuance
    \item \textbf{Users Service}: Managing buyer/seller profiles
    \item \textbf{Gig Service}: Managing jobs/services (gigs) with Elasticsearch search
    \item \textbf{Order Service}: Processing orders with Stripe payment integration
    \item \textbf{Chat Service}: Real-time communication via Socket.io
    \item \textbf{Review Service}: Rating and feedback for completed orders
    \item \textbf{Notification Service}: Sending emails and notifications (broker-driven)
\end{itemize}

Technologies used include \textbf{Node.js} and \textbf{TypeScript} for backend, \textbf{React/Vite} for frontend, \textbf{MongoDB/MySQL/PostgreSQL} for databases, \textbf{Redis} for caching, \textbf{RabbitMQ} for message broker, \textbf{Elasticsearch} for search, and \textbf{Kubernetes (K3s)} for deployment.

\cleardoublepage

% ---- MỤC LỤC ----
\tableofcontents
\cleardoublepage

% ---- DANH MỤC HÌNH VẼ ----
\listoffigures
\addcontentsline{toc}{chapter}{DANH MỤC HÌNH VẼ}
\cleardoublepage

% ---- DANH MỤC BẢNG ----
\listoftables
\addcontentsline{toc}{chapter}{DANH MỤC BẢNG}
\cleardoublepage

% ============================================================
\chapter{GIỚI THIỆU ĐỀ TÀI}
% ============================================================

\section{Đặt vấn đề}

Trong thời đại số hóa hiện nay, các nền tảng trực tuyến kết nối người cung cấp dịch vụ với người có nhu cầu đang phát triển mạnh mẽ. Các nền tảng như Fiverr, Upwork, Freelancer đã chứng minh mô hình ``gig economy'' (kinh tế giao việc) có tiềm năng to lớn. Tuy nhiên, việc xây dựng một hệ thống như vậy với kiến trúc nguyên khối (monolithic) sẽ gặp nhiều thách thức khi hệ thống lớn dần.

\textbf{Kiến trúc nguyên khối} là kiến trúc mà tất cả các thành phần của ứng dụng được gói gọn trong một dự án duy nhất. Mặc dù dễ xây dựng ban đầu, nhưng khi số lượng người dùng tăng lên, ứng dụng cần nhiều tính năng mới, dữ liệu tăng, logic phức tạp hơn --- ứng dụng sẽ trở thành ``con quái vật'' khó kiểm soát.

\textbf{Nhược điểm của kiến trúc nguyên khối:}
\begin{itemize}
    \item \textbf{Khó mở rộng}: Toàn bộ ứng dụng phải được scale cùng nhau, dù chỉ một chức năng nhỏ cần thêm tài nguyên
    \item \textbf{Thời gian triển khai lâu}: Mỗi thay đổi nhỏ đều yêu cầu build và deploy toàn bộ hệ thống
    \item \textbf{Rủi ro cao}: Một lỗi nhỏ có thể làm sập toàn bộ hệ thống
    \item \textbf{Khó cập nhật công nghệ}: Bị ràng buộc vào một ngôn ngữ/framework duy nhất
    \item \textbf{Khó bảo trì}: Codebase lớn khiến việc hiểu và sửa đổi trở nên khó khăn
\end{itemize}

\textbf{Bài toán đặt ra:} Cần tìm giải pháp kiến trúc cho phép:
\begin{itemize}
    \item Hệ thống có thể mở rộng từng phần độc lập (scale riêng từng chức năng)
    \item Triển khai nhanh, không làm gián đoạn dịch vụ (zero-downtime deployment)
    \item Tích hợp nhiều công nghệ phù hợp cho từng bài toán
    \item Đảm bảo độ sẵn sàng cao (high availability)
    \item Dễ dàng bảo trì và phát triển tính năng mới
\end{itemize}

\section{Mục tiêu và phạm vi đề tài}

\subsection{Mục tiêu}

\begin{itemize}
    \item \textbf{Mục tiêu 1}: Tìm hiểu sâu về kiến trúc Microservice, so sánh ưu nhược điểm với kiến trúc nguyên khối
    \item \textbf{Mục tiêu 2}: Thiết kế và xây dựng nền tảng trao đổi dịch vụ/công việc (gigs) theo mô hình Microservice với đầy đủ tính năng: xác thực, tìm kiếm, nhắn tin, thanh toán, đánh giá
    \item \textbf{Mục tiêu 3}: Áp dụng các công nghệ cloud-native hiện đại: Kubernetes, Docker, CI/CD, Message Broker, Elasticsearch
    \item \textbf{Mục tiêu 4}: Triển khai hệ thống trên môi trường thực tế (VPS + K3s) với khả năng tự động hóa
\end{itemize}

\subsection{Phạm vi đề tài}

\begin{itemize}
    \item \textbf{Phạm vi chức năng}: Hệ thống bao gồm 8 microservices xử lý các nghiệp vụ chính: Gateway, Authentication, User Management, Gig Management, Order Processing, Real-time Chat, Review System, Notification
    \item \textbf{Phạm vi công nghệ}: Node.js/TypeScript (Backend), React/Vite (Frontend), MongoDB/MySQL/PostgreSQL (Database), Redis (Cache), RabbitMQ (Message Broker), Elasticsearch (Search), Kubernetes/K3s (Orchestration)
    \item \textbf{Phạm vi triển khai}: Hệ thống được container hóa bằng Docker và triển khai trên Kubernetes với CI/CD tự động qua GitHub Actions
\end{itemize}

\section{Định hướng giải pháp}

\subsection{Lựa chọn kiến trúc Microservice}

Kiến trúc Microservice chia nhỏ ứng dụng thành các dịch vụ nhỏ, độc lập:
\begin{itemize}
    \item Mỗi service có thể phát triển, triển khai, scale độc lập
    \item Mỗi service có thể sử dụng công nghệ phù hợp nhất cho bài toán của mình (polyglot persistence, polyglot programming)
    \item Services giao tiếp qua lightweight protocols (HTTP/REST, message queue)
    \item Dễ dàng áp dụng CI/CD cho từng service riêng lẻ
\end{itemize}

\subsection{Thiết kế hệ thống ITHUST App}

Hệ thống được thiết kế gồm các thành phần chính:

\begin{itemize}
    \item \textbf{Frontend}: Ứng dụng React/Vite cung cấp giao diện người dùng hiện đại, responsive
    \item \textbf{API Gateway}: Điểm vào duy nhất cho client, xử lý routing, authentication, rate limiting
    \item \textbf{8 Microservices}: 
    \begin{itemize}
        \item Auth Service: JWT authentication, user registration/login
        \item Users Service: Profile management cho buyers và sellers
        \item Gig Service: CRUD gigs, Elasticsearch integration cho tìm kiếm
        \item Order Service: Order lifecycle, Stripe payment integration
        \item Chat Service: Real-time messaging qua Socket.io
        \item Review Service: Rating và feedback system
        \item Notification Service: Email notifications qua RabbitMQ
        \item Gateway Service: API Gateway
    \end{itemize}
    \item \textbf{Infrastructure}: 
    \begin{itemize}
        \item Message Broker (RabbitMQ): Giao tiếp bất đồng bộ giữa services
        \item Cache (Redis): Session storage, real-time adapter cho Socket.io
        \item Databases: MongoDB (document), MySQL (relational auth), PostgreSQL (reviews)
        \item Search Engine: Elasticsearch cho full-text search gigs
        \item Monitoring: Elastic APM, Kibana, Metricbeat
    \end{itemize}
\end{itemize}

\subsection{Công nghệ sử dụng}

\begin{itemize}
    \item \textbf{Backend}: Node.js 20+, Express.js, TypeScript
    \item \textbf{Frontend}: React 18, Vite, TypeScript, TailwindCSS, Redux Toolkit
    \item \textbf{Databases}: MongoDB (Mongoose), MySQL (Sequelize), PostgreSQL
    \item \textbf{Message Queue}: RabbitMQ (amqplib)
    \item \textbf{Real-time}: Socket.io với Redis Adapter
    \item \textbf{Search}: Elasticsearch 8.x (@elastic/elasticsearch)
    \item \textbf{Payment}: Stripe (@stripe/stripe-js)
    \item \textbf{Containerization}: Docker, Docker Compose
    \item \textbf{Orchestration}: Kubernetes (K3s - lightweight K8s)
    \item \textbf{CI/CD}: GitHub Actions
    \item \textbf{Monitoring}: Elastic APM, Pino (logging)
\end{itemize}

\section{Bố cục đồ án}

\begin{itemize}
    \item \textbf{Chương 1}: Giới thiệu đề tài --- Trình bày bối cảnh, vấn đề, mục tiêu và định hướng giải pháp
    \item \textbf{Chương 2}: Khảo sát và phân tích yêu cầu --- Phân tích ưu nhược điểm Microservice, đặc tả chức năng qua use case
    \item \textbf{Chương 3}: Công nghệ sử dụng --- Giới thiệu chi tiết các công nghệ: Node.js, React, MongoDB, RabbitMQ, Elasticsearch, Kubernetes
    \item \textbf{Chương 4}: Phát triển và triển khai --- Thiết kế kiến trúc, chi tiết 8 microservices, giao diện, triển khai K3s, CI/CD
    \item \textbf{Chương 5}: Kết luận và hướng phát triển --- Đánh giá kết quả và đề xuất cải tiến
\end{itemize}

% ============================================================
\chapter{KHẢO SÁT VÀ PHÂN TÍCH YÊU CẦU}
% ============================================================

\section{Khảo sát hiện trạng}

\subsection{Kiến trúc Microservice trong thực tế}

Kiến trúc Microservice đã được nhiều công ty công nghệ lớn áp dụng thành công:

\begin{itemize}
    \item \textbf{Netflix}: Hệ thống streaming với hàng trăm microservices, xử lý hàng tỷ request mỗi ngày
    \item \textbf{Amazon}: Kiến trúc service-oriented cho phép các đội nhóm nhỏ (two-pizza teams) phát triển độc lập
    \item \textbf{Uber}: Chuyển từ monolith sang microservices để hỗ trợ scale toàn cầu
    \item \textbf{Airbnb}: Microservices cho phép độc lập triển khai hàng trăm lần mỗi ngày
\end{itemize}

\subsection{Ưu điểm của kiến trúc Microservice}

\begin{itemize}
    \item \textbf{Scale độc lập}: Chỉ scale các service cần thiết, tiết kiệm tài nguyên
    \item \textbf{Triển khai độc lập}: Mỗi service có thể deploy riêng, giảm rủi ro, hỗ trợ zero-downtime
    \item \textbf{Công nghệ đa dạng}: Mỗi service chọn công nghệ phù hợp nhất (polyglot)
    \item \textbf{Fault isolation}: Lỗi một service không làm sập toàn hệ thống
    \item \textbf{Tổ chức đội nhóm}: Các đội nhỏ (5-10 người) phụ trách từng service, tăng năng suất
    \item \textbf{Dễ bảo trì}: Codebase nhỏ, dễ hiểu và thay đổi
\end{itemize}

\subsection{Nhược điểm và thách thức}

\begin{itemize}
    \item \textbf{Độ phức tạp phân tán}: Phải xử lý network latency, partial failures, distributed transactions
    \item \textbf{Operational overhead}: Cần công cụ quản lý, monitoring, logging tập trung
    \item \textbf{Data consistency}: Dữ liệu phân tán nhiều nơi, khó đảm bảo tính nhất quán toàn cục
    \item \textbf{Testing phức tạp}: Cần integration testing, contract testing giữa services
    \item \textbf{DevOps yêu cầu cao}: Cần CI/CD, container orchestration (Kubernetes)
    \item \textbf{Theo nguyên tắc CAP}: Phải đánh đổi giữa Consistency, Availability, Partition Tolerance
\end{itemize}

\subsection{Khi nào sử dụng Microservice?}

Kiến trúc Microservice phù hợp khi:
\begin{itemize}
    \item Hệ thống lớn, nhiều đội phát triển
    \item Cần scale từng phần riêng lẻ
    \item Yêu cầu độ sẵn sàng cao (high availability)
    \item Cần triển khai liên tục (continuous deployment)
    \item Các domain có thể tách biệt rõ ràng
\end{itemize}

Không nên dùng khi:
\begin{itemize}
    \item Đội nhóm nhỏ (dưới 10 người)
    \item Ứng dụng đơn giản, ít chức năng
    \item Chưa có kinh nghiệm DevOps
    \item Không có yêu cầu scale phức tạp
\end{itemize}

\section{Tổng quan chức năng}

\subsection{Biểu đồ use case tổng quan}

\begin{figure}[H]
    \centering
    \fbox{\parbox{0.8\textwidth}{
        \centering
        \textbf{[Đề xuất ảnh]}\\
        Biểu đồ Use Case tổng quan ITHUST App\\[0.5cm]
        - Actors: Guest, Buyer, Seller, Admin\\
        - Use Cases: Đăng ký/Đăng nhập, Tìm kiếm Gigs, Đặt hàng, Nhắn tin, Đánh giá, Quản lý Profile
    }}
    \caption{Biểu đồ use case tổng quan}
    \label{fig:uc_tongquan}
\end{figure}

Hệ thống có 4 tác nhân chính:
\begin{itemize}
    \item \textbf{Guest}: Khách chưa đăng nhập, có thể tìm kiếm và xem gigs
    \item \textbf{Buyer}: Người mua dịch vụ, có thể đặt hàng, thanh toán, chat với seller
    \item \textbf{Seller}: Người bán dịch vụ, tạo và quản lý gigs, xử lý đơn hàng
    \item \textbf{Admin}: Quản trị viên, quản lý người dùng và nội dung
\end{itemize}

\subsection{Biểu đồ use case phân rã --- Xác thực}

\begin{figure}[H]
    \centering
    \fbox{\parbox{0.8\textwidth}{
        \centering
        \textbf{[Đề xuất ảnh]}\\
        Use Case: Đăng ký, Đăng nhập, Quên mật khẩu, Xác thực email\\
        Đăng xuất, Đổi mật khẩu, Cập nhật profile
    }}
    \caption{Biểu đồ use case phân rã chức năng xác thực}
    \label{fig:uc_auth}
\end{figure}

\subsection{Biểu đồ use case phân rã --- Quản lý Gigs}

\begin{figure}[H]
    \centering
    \fbox{\parbox{0.8\textwidth}{
        \centering
        \textbf{[Đề xuất ảnh]}\\
        Use Case (Seller): Tạo gig, Sửa gig, Xóa gig, Xem danh sách gigs\\
        Use Case (Buyer): Tìm kiếm gigs (Elasticsearch), Xem chi tiết gig, So sánh gigs
    }}
    \caption{Biểu đồ use case phân rã chức năng quản lý Gigs}
    \label{fig:uc_gig}
\end{figure}

\subsection{Biểu đồ use case phân rã --- Đặt hàng và Thanh toán}

\begin{figure}[H]
    \centering
    \fbox{\parbox{0.8\textwidth}{
        \centering
        \textbf{[Đề xuất ảnh]}\\
        Use Case: Thêm vào giỏ, Tạo đơn hàng, Thanh toán (Stripe),\\
        Xem lịch sử đơn hàng, Hủy đơn, Theo dõi trạng thái
    }}
    \caption{Biểu đồ use case phân rã chức năng đặt hàng}
    \label{fig:uc_order}
\end{figure}

\subsection{Biểu đồ use case phân rã --- Chat và Đánh giá}

\begin{figure}[H]
    \centering
    \fbox{\parbox{0.8\textwidth}{
        \centering
        \textbf{[Đề xuất ảnh]}\\
        Chat: Gửi tin nhắn (real-time), Nhận thông báo tin nhắn mới,\\
        Xem lịch sử chat\\[0.3cm]
        Review: Đánh giá đơn hàng, Xem đánh giá của seller, Phản hồi đánh giá
    }}
    \caption{Biểu đồ use case phân rã chức năng Chat và Review}
    \label{fig:uc_chat_review}
\end{figure}

\section{Đặc tả chức năng}

\subsection{Đặc tả use case --- Tìm kiếm Gigs với Elasticsearch}

\begin{table}[H]
\centering
\begin{tabular}{|p{4cm}|p{10cm}|}
\hline
\textbf{Tác nhân} & Buyer/Guest \\ \hline
\multirow{6}{*}{\textbf{Luồng chính}} &
1. Người dùng nhập từ khóa tìm kiếm \\ \cline{2-2}
& 2. Hệ thống gửi query đến Elasticsearch \\ \cline{2-2}
& 3. Elasticsearch trả về kết quả matching (fuzzy search, relevance scoring) \\ \cline{2-2}
& 4. Hệ thống hiển thị danh sách gigs với thông tin: tiêu đề, giá, rating, seller \\ \cline{2-2}
& 5. Người dùng có thể lọc theo: danh mục, giá, delivery time, seller level \\ \hline
\multirow{2}{*}{\textbf{Luồng thay thế}} &
1. Không có kết quả: Hiển thị ``Không tìm thấy kết quả'' + gợi ý từ khóa \\ \cline{2-2}
& 2. Elasticsearch không khả dụng: Fallback về tìm kiếm cơ bản từ MongoDB \\ \hline
\end{tabular}
\caption{Đặc tả use case tìm kiếm Gigs}
\label{tab:uc_search}
\end{table}

\subsection{Đặc tả use case --- Đặt hàng và Thanh toán}

\begin{table}[H]
\centering
\begin{tabular}{|p{4cm}|p{10cm}|}
\hline
\textbf{Tác nhân} & Buyer \\ \hline
\multirow{8}{*}{\textbf{Luồng chính}} &
1. Buyer chọn gig và click ``Order Now'' \\ \cline{2-2}
& 2. Hệ thống tạo order với trạng thái ``pending'' \\ \cline{2-2}
& 3. Buyer chọn phương thức thanh toán (Stripe Card/Wallet) \\ \cline{2-2}
& 4. Hệ thống gọi Stripe API để xử lý thanh toán \\ \cline{2-2}
& 5. Thanh toán thành công: Order chuyển sang ``in progress'' \\ \cline{2-2}
& 6. Notification Service gửi email cho Seller qua RabbitMQ \\ \cline{2-2}
& 7. Buyer và Seller có thể chat để trao đổi yêu cầu \\ \cline{2-2}
& 8. Seller hoàn thành và giao hàng, Buyer xác nhận \\ \hline
\multirow{2}{*}{\textbf{Luồng thay thế}} &
4a. Thanh toán thất bại: Hiển thị lỗi, cho phép thử lại hoặc đổi phương thức \\ \cline{2-2}
& 8a. Buyer không xác nhận: Tự động hoàn thành sau 3 ngày \\ \hline
\end{tabular}
\caption{Đặc tả use case đặt hàng và thanh toán}
\label{tab:uc_order}
\end{table}

\subsection{Đặc tả use case --- Chat Real-time}

\begin{table}[H]
\centering
\begin{tabular}{|p{4cm}|p{10cm}|}
\hline
\textbf{Tác nhân} & Buyer, Seller \\ \hline
\multirow{6}{*}{\textbf{Luồng chính}} &
1. User mở conversation với user khác \\ \cline{2-2}
& 2. Frontend kết nối Socket.io đến Chat Service \\ \cline{2-2}
& 3. Chat Service lưu message vào MongoDB, broadcast qua Redis Adapter \\ \cline{2-2}
& 4. Receiver nhận message real-time qua Socket.io \\ \cline{2-2}
& 5. Hiển thị trạng thái ``đã đọc'' khi receiver mở conversation \\ \cline{2-2}
& 6. Lưu trữ lịch sử chat có thể truy xuất sau \\ \hline
\textbf{Luồng thay thế} &
4a. Receiver offline: Message được lưu, gửi push notification qua Notification Service \\ \hline
\end{tabular}
\caption{Đặc tả use case chat real-time}
\label{tab:uc_chat}
\end{table}

% ============================================================
\chapter{CÔNG NGHỆ SỬ DỤNG}
% ============================================================

\section{Node.js và TypeScript}

\subsection{Node.js}

\textbf{Node.js} là môi trường runtime JavaScript built on Chrome's V8 JavaScript engine, cho phép chạy JavaScript ở server-side.

\textbf{Đặc điểm chính:}
\begin{itemize}
    \item \textbf{Non-blocking I/O}: Xử lý nhiều request đồng thời mà không chờ đợi (event-driven, asynchronous)
    \item \textbf{Single-threaded with Event Loop}: Hiệu quả cho I/O-bound operations
    \item \textbf{NPM ecosystem}: Hơn 1 triệu packages sẵn có
    \item \textbf{Cross-platform}: Chạy trên Windows, Linux, macOS
\end{itemize}

\textbf{Tại sao chọn Node.js cho Microservice:}
\begin{itemize}
    \item Phù hợp cho I/O-intensive services (API Gateway, Chat, Notification)
    \item Khởi động nhanh, phù hợp containerization
    \item JSON-native, dễ dàng xây dựng REST APIs
    \item Cộng đồng lớn, nhiều thư viện hỗ trợ microservices
\end{itemize}

\subsection{TypeScript}

\textbf{TypeScript} là superset của JavaScript, thêm static typing vào JS.

\textbf{Lợi ích:}
\begin{itemize}
    \item Type safety: Phát hiện lỗi tại compile-time thay vì runtime
    \item IntelliSense: Code completion, refactoring tốt hơn trong IDE
    \item Documentation tự động: Types như tài liệu sống
    \item OOP features: Interfaces, Generics, Enums, Decorators
    \item Compile to ES5/ES6: Tương thích các trình duyệt
\end{itemize}

\begin{lstlisting}[style=TypeScriptStyle, caption={Ví dụ TypeScript Interface}]
interface IGig {
  _id: string;
  sellerId: string;
  title: string;
  description: string;
  categories: string;
  subCategories: string[];
  tags: string[];
  active: boolean;
  expectedDelivery: string;
  price: number;
  sortId: number;
}

// Type safety ensures correct usage
const createGig = (gigData: IGig): Promise<IGig> => {
  return GigModel.create(gigData);
};
\end{lstlisting}

\section{React và Vite}

\subsection{React 18}

\textbf{React} là thư viện JavaScript để xây dựng user interfaces, phát triển bởi Facebook.

\textbf{Đặc điểm chính:}
\begin{itemize}
    \item \textbf{Component-based}: UI được chia thành các components độc lập, reusable
    \item \textbf{Virtual DOM}: Cập nhật UI hiệu quả, chỉ render những gì thay đổi
    \item \textbf{Unidirectional data flow}: Dữ liệu chảy một chiều, dễ debug
    \item \textbf{Hooks}: useState, useEffect, useContext cho functional components
\end{itemize}

\subsection{Vite}

\textbf{Vite} là next generation frontend build tool, thay thế Webpack.

\textbf{Ưu điểm:}
\begin{itemize}
    \item \textbf{Dev server cực nhanh}: Sử dụng native ES modules, không cần bundling khi dev
    \item \textbf{Hot Module Replacement (HMR)}: Cập nhật tức thì không reload trang
    \item \textbf{Build tối ưu}: Rollup-based production build với tree-shaking
    \item \textbf{TypeScript support built-in}: Không cần cấu hình phức tạp
\end{itemize}

\begin{lstlisting}[style=TypeScriptStyle, caption={Vite Configuration}]
// vite.config.ts
import { defineConfig } from 'vite';
import react from '@vitejs/plugin-react';
import tsconfigPaths from 'vite-tsconfig-paths';

export default defineConfig({
  plugins: [react(), tsconfigPaths()],
  server: {
    port: 3000,
    proxy: {
      '/api': 'http://localhost:4000'
    }
  },
  build: {
    outDir: 'build',
    sourcemap: true
  }
});
\end{lstlisting}

\section{Cơ sở dữ liệu}

\subsection{MongoDB}

\textbf{MongoDB} là document database NoSQL, lưu trữ dữ liệu dạng JSON-like (BSON).

\textbf{Sử dụng cho:} Gig Service, Chat Service (lưu messages), Notification Service

\textbf{Lý do chọn:}
\begin{itemize}
    \item Schema flexible phù hợp dữ liệu gigs có cấu trúc đa dạng
    \item Horizontal scaling qua sharding
    \item Rich query language với aggregation pipeline
    \item Native JSON, tương thích với JavaScript
\end{itemize}

\subsection{MySQL}

\textbf{MySQL} là relational database, sử dụng cho Auth Service.

\textbf{Lý do chọn:}
\begin{itemize}
    \item Dữ liệu user authentication có cấu trúc rõ ràng, cần ACID compliance
    \item Sequelize ORM hỗ trợ migrations, transactions
    \item Phù hợp cho relational data: users, roles, permissions
\end{itemize}

\subsection{PostgreSQL}

\textbf{PostgreSQL} là advanced relational database, sử dụng cho Review Service.

\textbf{Lý do chọn:}
\begin{itemize}
    \item Hỗ trợ complex queries cho thống kê đánh giá
    \item JSON/JSONB support cho semi-structured data
    \item Robust và reliable cho critical data
\end{itemize}

\section{Redis}

\textbf{Redis} (REmote DIctionary Server) là in-memory data structure store, dùng làm database, cache, message broker.

\textbf{Sử dụng trong hệ thống:}
\begin{itemize}
    \item \textbf{Session storage}: Lưu JWT sessions tạm thời
    \item \textbf{Socket.io Adapter}: Cho phép multiple Chat Service instances scale horizontally
    \item \textbf{Rate limiting}: Track request counts trong window time
    \item \textbf{Caching}: Cache frequently accessed data
\end{itemize}

\textbf{Redis Adapter cho Socket.io:}
\begin{lstlisting}[style=TypeScriptStyle, caption={Redis Adapter Configuration}]
import { createAdapter } from '@socket.io/redis-adapter';
import { createClient } from 'redis';

const pubClient = createClient({ url: redisUrl });
const subClient = pubClient.duplicate();

io.adapter(createAdapter(pubClient, subClient));
// Now multiple server instances can broadcast to all clients
\end{lstlisting}

\section{RabbitMQ}

\subsection{Giới thiệu}

\textbf{RabbitMQ} là message broker mã nguồn mở, implement giao thức AMQP (Advanced Message Queuing Protocol).

\textbf{Khái niệm chính:}
\begin{itemize}
    \item \textbf{Producer}: Ứng dụng gửi message
    \item \textbf{Exchange}: Nhận message từ producer, route đến queue phù hợp
    \item \textbf{Queue}: Lưu trữ message cho đến khi consumer xử lý
    \item \textbf{Consumer}: Ứng dụng nhận và xử lý message
\end{itemize}

\subsection{Exchange Types}

\begin{itemize}
    \item \textbf{Direct}: Route message đến queue có routing key khớp chính xác
    \item \textbf{Fanout}: Broadcast đến tất cả queues bound
    \item \textbf{Topic}: Route theo pattern matching (wildcards: * và #)
    \item \textbf{Headers}: Route dựa trên headers của message
\end{itemize}

\subsection{Sử dụng trong ITHUST App}

\begin{lstlisting}[style=TypeScriptStyle, caption={RabbitMQ Message Publishing}]
// Order Service publishes notification
import { Channel } from 'amqplib';

const publishOrderNotification = async (
  channel: Channel,
  orderData: IOrderMessage
): Promise<void> => {
  const exchangeName = 'order-events';
  const routingKey = 'order.created';
  
  await channel.assertExchange(exchangeName, 'topic', { durable: true });
  
  channel.publish(
    exchangeName,
    routingKey,
    Buffer.from(JSON.stringify(orderData)),
    { persistent: true } // Message survives broker restart
  );
};
\end{lstlisting}

\section{Elasticsearch}

\subsection{Giới thiệu}

\textbf{Elasticsearch} là distributed search and analytics engine, built on Apache Lucene.

\textbf{Đặc điểm:}
\begin{itemize}
    \item \textbf{Full-text search}: Powerful text analysis, tokenization, relevance scoring
    \item \textbf{Near real-time}: Search results available trong vòng 1 giây
    \item \textbf{Horizontal scaling}: Thêm nodes để tăng capacity
    \item \textbf{Schema-free}: JSON documents, dynamic mapping
\end{itemize}

\subsection{Integration với Node.js}

\begin{lstlisting}[style=TypeScriptStyle, caption={Elasticsearch Client}]
import { Client } from '@elastic/elasticsearch';

const elasticClient = new Client({
  node: process.env.ELASTICSEARCH_URL,
  auth: {
    username: process.env.ELASTIC_USER!,
    password: process.env.ELASTIC_PASSWORD!
  }
});

// Index a gig
await elasticClient.index({
  index: 'gigs',
  id: gig._id,
  document: {
    title: gig.title,
    description: gig.description,
    categories: gig.categories,
    tags: gig.tags,
    price: gig.price,
    sellerId: gig.sellerId
  }
});

// Search with relevance scoring
const result = await elasticClient.search({
  index: 'gigs',
  query: {
    multi_match: {
      query: searchTerm,
      fields: ['title^3', 'description', 'tags^2'],
      fuzziness: 'AUTO'
    }
  }
});
\end{lstlisting}

\section{Socket.io}

\textbf{Socket.io} là library cho real-time, bidirectional, event-based communication.

\textbf{Features:}
\begin{itemize}
    \item \textbf{WebSocket with fallbacks}: WebSocket, AJAX long-polling, etc.
    \item \textbf{Rooms and Namespaces}: Tổ chức connections
    \item \textbf{Automatic reconnection}: Client tự động reconnect khi mất kết nối
    \item \textbf{Binary support}: Hỗ trợ ArrayBuffer, Blob
\end{itemize}

\begin{lstlisting}[style=TypeScriptStyle, caption={Socket.io Chat Implementation}]
// Server side (Chat Service)
io.on('connection', (socket) => {
  console.log(`User connected: ${socket.id}`);
  
  socket.on('join-conversation', (conversationId) => {
    socket.join(conversationId);
  });
  
  socket.on('send-message', async (data) => {
    const message = await saveMessageToDB(data);
    
    // Broadcast to all clients in the conversation
    io.to(data.conversationId).emit('new-message', message);
    
    // Notify offline users via RabbitMQ
    await publishNotification(data.receiverId, 'new-message');
  });
  
  socket.on('disconnect', () => {
    console.log(`User disconnected: ${socket.id}`);
  });
});
\end{lstlisting}

\section{Stripe Payment}

\textbf{Stripe} là payment processing platform, cung cấp APIs cho thanh toán online.

\textbf{Tích hợp trong Order Service:}
\begin{itemize}
    \item \textbf{Payment Intent}: Tạo và xác nhận thanh toán
    \item \textbf{Webhooks}: Nhận events từ Stripe (payment success, failure)
    \item \textbf{PCI Compliance}: Stripe Elements xử lý card data an toàn
\end{itemize}

\section{Docker và Kubernetes (K3s)}

\subsection{Docker}

\textbf{Docker} là platform để develop, ship, run applications trong containers.

\textbf{Dockerfile cho Node.js service:}
\begin{lstlisting}[style=BashStyle, caption={Dockerfile cho Gateway Service}]
# Build stage
FROM node:20-alpine AS builder
WORKDIR /app
COPY package*.json ./
RUN npm ci
COPY . .
RUN npm run build

# Production stage
FROM node:20-alpine
WORKDIR /app
COPY --from=builder /app/build ./build
COPY --from=builder /app/node_modules ./node_modules
COPY package*.json ./
EXPOSE 4000
CMD ["node", "build/src/app.js"]
\end{lstlisting}

\subsection{Kubernetes và K3s}

\textbf{Kubernetes} là container orchestration platform tự động hóa deployment, scaling, management.

\textbf{K3s} là lightweight Kubernetes distribution, phù hợp cho:
\begin{itemize}
    \item Edge computing, IoT
    \item CI/CD pipelines
    \item Single-node và small-scale deployments (ARM, VPS)
\end{itemize}

\textbf{K3s Architecture:}
\begin{itemize}
    \item Single binary dưới 100MB
    \item Built-in Traefik Ingress Controller
    \item SQLite thay vì etcd (default)
    \item Containerd runtime (thay vì Docker)
\end{itemize}

\section{CI/CD với GitHub Actions}

\textbf{GitHub Actions} là platform tự động hóa CI/CD trực tiếp trong GitHub repository.

\textbf{Workflow cho ITHUST App:}
\begin{enumerate}
    \item \textbf{Trigger}: Push/PR vào main branch
    \item \textbf{Path filtering}: Chỉ build services có thay đổi
    \item \textbf{Build}: Docker build với tag \texttt{stable-<build-number>}
    \item \textbf{Push}: Push image lên Docker Hub
    \item \textbf{Deploy}: Kubectl apply rolling update đến K3s cluster
\end{enumerate}

\begin{lstlisting}[style=BashStyle, caption={GitHub Actions Workflow (đơn giản hóa)}]
name: CI/CD Pipeline

on:
  push:
    branches: [ main ]
    paths:
      - 'server/1-gateway-service/**'

jobs:
  build-and-deploy:
    runs-on: ubuntu-latest
    steps:
      - uses: actions/checkout@v4
      
      - name: Docker Login
        run: echo ${{ secrets.DOCKER_PASSWORD }} | docker login -u ${{ secrets.DOCKER_USERNAME }} --password-stdin
      
      - name: Docker Build
        run: |
          docker build -t ${{ secrets.DOCKER_USERNAME }}/gateway:stable-${{ github.run_number }} \
            ./server/1-gateway-service
      
      - name: Docker Push
        run: docker push ${{ secrets.DOCKER_USERNAME }}/gateway:stable-${{ github.run_number }}
      
      - name: Deploy to K3s
        run: |
          echo "${{ secrets.KUBECONFIG_B64 }}" | base64 -d > kubeconfig
          kubectl --kubeconfig=kubeconfig set image deployment/gateway \
            gateway=${{ secrets.DOCKER_USERNAME }}/gateway:stable-${{ github.run_number }}
\end{lstlisting}

% ============================================================
\chapter{PHÁT TRIỂN VÀ TRIỂN KHAI ỨNG DỤNG}
% ============================================================

\section{Thiết kế kiến trúc}

\subsection{Kiến trúc tổng quan}

\begin{figure}[H]
    \centering
    \fbox{\parbox{0.9\textwidth}{
        \centering
        \textbf{[Đề xuất ảnh cần chụp]}\\[0.3cm]
        \textbf{Kiến trúc tổng quan ITHUST App}\\[0.3cm]
        - Frontend (React/Vite) $\rightarrow$ API Gateway\\
        - API Gateway $\rightarrow$ 8 Microservices\\
        - Services $\leftrightarrow$ Message Broker (RabbitMQ)\\
        - Services $\leftrightarrow$ Databases (MongoDB, MySQL, PostgreSQL)\\
        - Services $\leftrightarrow$ Cache (Redis)\\
        - Gig Service $\leftrightarrow$ Elasticsearch\\
        - Chat Service $\leftrightarrow$ Socket.io + Redis Adapter\\
        - Order Service $\leftrightarrow$ Stripe API\\
        - Kubernetes (K3s) orchestration
    }}
    \caption{Kiến trúc tổng quan hệ thống}
    \label{fig:kientruc_tongquan}
\end{figure}

\subsection{API Gateway Pattern}

\textbf{Vai trò của Gateway Service:}
\begin{itemize}
    \item \textbf{Single entry point}: Tất cả client requests đi qua gateway
    \item \textbf{Routing}: Chuyển hướng đến service phù hợp dựa trên path
    \item \textbf{Authentication}: Verify JWT token trước khi forward request
    \item \textbf{Rate limiting}: Giới hạn request/minute để chống DDoS
    \item \textbf{Load balancing}: Phân phối request đến multiple instances
    \item \textbf{Response aggregation}: Kết hợp responses từ nhiều services
\end{itemize}

\begin{lstlisting}[style=TypeScriptStyle, caption={Gateway Routing}]
// Gateway routes configuration
app.use('/api/v1/auth', proxy('http://auth-service:4001'));
app.use('/api/v1/users', proxy('http://users-service:4002'));
app.use('/api/v1/gigs', proxy('http://gig-service:4003'));
app.use('/api/v1/orders', proxy('http://order-service:4004'));
app.use('/api/v1/chat', proxy('http://chat-service:4005'));
app.use('/api/v1/reviews', proxy('http://review-service:4006'));
app.use('/api/v1/notifications', proxy('http://notification-service:4007'));
\end{lstlisting}

\subsection{Service Communication Patterns}

\textbf{Synchronous (REST API):} Client $\rightarrow$ Gateway $\rightarrow$ Service $\rightarrow$ Database
\begin{itemize}
    \item Ưu điểm: Đơn giản, response ngay lập tức
    \item Nhược điểm: Tight coupling, cascade failures
    \item Dùng cho: Read operations, request-response pattern
\end{itemize}

\textbf{Asynchronous (Message Queue):} Service A $\rightarrow$ RabbitMQ $\rightarrow$ Service B
\begin{itemize}
    \item Ưu điểm: Loose coupling, fault tolerance, buffering
    \item Nhược điểm: Complexity, eventual consistency
    \item Dùng cho: Notifications, background jobs, event propagation
\end{itemize}

\subsection{Thiết kế cơ sở dữ liệu theo service}

\textbf{Database-per-Service Pattern:} Mỗi service có database riêng, đảm bảo:
\begin{itemize}
    \item Loose coupling giữa services
    \item Độc lập phát triển và triển khai
    \item Polyglot persistence (chọn DB phù hợp cho từng bài toán)
\end{itemize}

\begin{table}[H]
\centering
\begin{tabular}{|l|l|l|}
\hline
\textbf{Service} & \textbf{Database} & \textbf{Lý do} \\ \hline
Auth Service & MySQL & Relational, ACID cho user data \\ \hline
Users Service & MongoDB & Flexible schema cho profiles \\ \hline
Gig Service & MongoDB + Elasticsearch & Document + Full-text search \\ \hline
Order Service & MongoDB & Order lifecycle, transactions \\ \hline
Chat Service & MongoDB & Message storage, time-series \\ \hline
Review Service & PostgreSQL & Complex analytics queries \\ \hline
Notification Service & MongoDB & Queue-like message storage \\ \hline
Gateway Service & Redis & Rate limiting, session cache \\ \hline
\end{tabular}
\caption{Database assignment cho từng service}
\label{tab:database_assignment}
\end{table}

\section{Thiết kế chi tiết các microservices}

\subsection{Auth Service}

\textbf{Chức năng:} Xác thực và phân quyền người dùng

\textbf{APIs:}
\begin{itemize}
    \item \texttt{POST /api/v1/auth/signup} --- Đăng ký tài khoản mới
    \item \texttt{POST /api/v1/auth/signin} --- Đăng nhập, trả về JWT
    \item \texttt{POST /api/v1/auth/verify-email} --- Xác thực email
    \item \texttt{POST /api/v1/auth/forgot-password} --- Quên mật khẩu
    \item \texttt{POST /api/v1/auth/reset-password} --- Đặt lại mật khẩu
    \item \texttt{GET /api/v1/auth/validate-token} --- Kiểm tra token hợp lệ
\end{itemize}

\textbf{Security:}
\begin{itemize}
    \item Bcrypt cho password hashing (salt rounds: 12)
    \item JWT với RS256 (asymmetric signing)
    \item Rate limiting: 5 attempts / 15 minutes cho auth endpoints
    \item CORS policy cho cross-origin requests
\end{itemize}

\subsection{Users Service}

\textbf{Chức năng:} Quản lý hồ sơ người dùng (buyers và sellers)

\textbf{APIs:}
\begin{itemize}
    \item \texttt{GET /api/v1/users/:userId} --- Lấy profile user
    \item \texttt{PATCH /api/v1/users/:userId} --- Cập nhật profile
    \item \texttt{GET /api/v1/users/seller/:sellerId} --- Lấy thông tin seller với stats
    \item \texttt{POST /api/v1/users/seller/create} --- Đăng ký làm seller
\end{itemize}

\textbf{Events published:}
\begin{itemize}
    \item \texttt{user.profile.updated} --- Thông báo các service khác cập nhật cache
\end{itemize}

\subsection{Gig Service}

\textbf{Chức năng:} Quản lý dịch vụ/công việc (gigs) với tìm kiếm full-text

\textbf{APIs:}
\begin{itemize}
    \item CRUD operations cho gigs
    \item \texttt{GET /api/v1/gigs/search/:searchTerm} --- Tìm kiếm Elasticsearch
    \item \texttt{GET /api/v1/gigs/category/:category} --- Lọc theo danh mục
    \item \texttt{GET /api/v1/gigs/seller/:sellerId} --- Lấy gigs của seller
\end{itemize}

\textbf{Elasticsearch Integration:}
\begin{lstlisting}[style=TypeScriptStyle, caption={Gig Search Implementation}]
export const searchGigs = async (searchTerm: string): Promise<ISearchResult> => {
  const result: SearchResponse<IGig> = await elasticClient.search({
    index: 'gigs',
    query: {
      bool: {
        must: [
          {
            multi_match: {
              query: searchTerm,
              fields: ['title^3', 'description', 'tags^2', 'categories'],
              type: 'best_fields',
              fuzziness: 'AUTO'
            }
          }
        ],
        filter: [
          { term: { active: true } }
        ]
      }
    },
    sort: [
      { sortId: 'desc' },
      { rating: 'desc' }
    ],
    size: 50
  });
  
  return {
    hits: result.hits.hits.map(hit => hit._source),
    total: result.hits.total
  };
};
\end{lstlisting}

\subsection{Order Service}

\textbf{Chức trạng thái đơn hàng:}
\begin{itemize}
    \item \texttt{pending} --- Chờ thanh toán
    \item \texttt{in-progress} --- Đang thực hiện
    \item \texttt{delivered} --- Đã giao, chờ xác nhận
    \item \texttt{completed} --- Hoàn thành
    \item \texttt{cancelled} --- Đã hủy
    \item \texttt{disputed} --- Có tranh chấp
\end{itemize}

\textbf{Payment Flow với Stripe:}
\begin{enumerate}
    \item Frontend tạo Payment Intent qua Stripe.js
    \item Stripe trả về client secret
    \item User nhập card info, Stripe Elements xử lý securely
    \item Stripe webhook gọi \texttt{/stripe-webhook} khi payment success
    \item Order chuyển từ \texttt{pending} $\rightarrow$ \texttt{in-progress}
    \item Publish message \texttt{order.created} đến Notification Service
\end{enumerate}

\subsection{Chat Service}

\textbf{Kiến trúc real-time:}
\begin{itemize}
    \item Socket.io server với Redis Adapter cho horizontal scaling
    \item MongoDB lưu trữ conversations và messages
    \item Message schema: \texttt{\{conversationId, senderId, receiverId, content, createdAt, read\}}
\end{itemize}

\textbf{Events:}
\begin{itemize}
    \item \texttt{connection} --- User kết nối
    \item \texttt{join-conversation} --- Tham gia room chat
    \item \texttt{send-message} --- Gửi tin nhắn
    \item \texttt{new-message} --- Nhận tin nhắn mới
    \item \texttt{typing} --- Trạng thái đang gõ
    \item \texttt{message-read} --- Đánh dấu đã đọc
\end{itemize}

\subsection{Review Service}

\textbf{Chức năng:} Đánh giá và phản hồi sau khi đơn hàng hoàn thành

\textbf{APIs:}
\begin{itemize}
    \item \texttt{POST /api/v1/reviews} --- Tạo đánh giá (chỉ sau khi order completed)
    \item \texttt{GET /api/v1/reviews/gig/:gigId} --- Lấy đánh giá của gig
    \item \texttt{GET /api/v1/reviews/seller/:sellerId} --- Lấy đánh giá của seller với stats
    \item \texttt{PATCH /api/v1/reviews/:reviewId} --- Seller phản hồi đánh giá
\end{itemize}

\subsection{Notification Service}

\textbf{Kiến trúc Event-Driven:} Consumer từ RabbitMQ, không có REST API public

\textbf{Xử lý events:}
\begin{itemize}
    \item \texttt{order.created} --- Gửi email cho seller
    \item \texttt{order.completed} --- Gửi email cho cả buyer và seller
    \item \texttt{message.received} --- Push notification cho offline users
    \item \texttt{user.welcome} --- Email chào mừng khi đăng ký
\end{itemize}

\begin{lstlisting}[style=TypeScriptStyle, caption={RabbitMQ Consumer}]
const consumeMessages = async (): Promise<void> => {
  const connection = await amqp.connect(rabbitmqUrl);
  const channel = await connection.createChannel();
  
  await channel.assertExchange('order-events', 'topic', { durable: true });
  const { queue } = await channel.assertQueue('notification-queue', { durable: true });
  
  await channel.bindQueue(queue, 'order-events', 'order.*');
  
  channel.consume(queue, async (msg) => {
    if (msg) {
      const event = JSON.parse(msg.content.toString());
      
      switch (event.type) {
        case 'order.created':
          await sendEmailToSeller(event.sellerId, 'New order received');
          break;
        case 'order.completed':
          await sendCompletionEmails(event.buyerId, event.sellerId);
          break;
      }
      
      channel.ack(msg); // Acknowledge message processed
    }
  });
};
\end{lstlisting}

\section{Thiết kế giao diện}

\subsection{Công nghệ UI}

\textbf{Stack:}
\begin{itemize}
    \item \textbf{React 18} với Functional Components và Hooks
    \item \textbf{TypeScript} cho type safety
    \item \textbf{Vite} cho dev server nhanh và optimized builds
    \item \textbf{TailwindCSS} cho utility-first styling
    \item \textbf{Redux Toolkit} cho state management
    \item \textbf{React Router v6} cho routing
\end{itemize}

\subsection{Layout và Components chính}

\begin{figure}[H]
    \centering
    \fbox{\parbox{0.85\textwidth}{
        \centering
        \textbf{[Đề xuất ảnh cần chụp]}\\[0.3cm]
        \textbf{Giao diện chính ITHUST App}\\[0.3cm]
        - Header: Logo, Search bar, Navigation, User menu\\
        - Hero section: Categories, Call-to-action\\
        - Featured Gigs: Grid layout với gig cards\\
        - Footer: Links, Social media
    }}
    \caption{Layout tổng quan ứng dụng}
    \label{fig:layout}
\end{figure}

\subsection{Các màn hình chính cần chụp}

\textbf{Danh sách ảnh cần chụp cho báo cáo:}

\begin{enumerate}
    \item \textbf{Trang chủ (Homepage)}: Hero section, featured gigs, categories
    \item \textbf{Trang tìm kiếm (Search Results)}: Filters, grid kết quả, pagination
    \item \textbf{Chi tiết Gig (Gig Detail)}: Gallery, mô tả, packages, reviews, seller info
    \item \textbf{Giỏ hàng/Checkout}: Order summary, Stripe payment form
    \item \textbf{Dashboard Seller}: Quản lý gigs, orders, earnings
    \item \textbf{Dashboard Buyer}: Orders history, messages
    \item \textbf{Chat Interface}: Conversation list, chat window, real-time messages
    \item \textbf{Profile/User Settings}: Avatar, thông tin cá nhân, đổi mật khẩu
    \item \textbf{Login/Register}: Forms với validation
\end{enumerate}

\section{Triển khai trên Kubernetes (K3s)}

\subsection{K8s Resources cho mỗi Service}

\textbf{Deployment.yaml} --- Định nghĩa pods và replicas:
\begin{lstlisting}[style=BashStyle, caption={K8s Deployment (đơn giản hóa)}]
apiVersion: apps/v1
kind: Deployment
metadata:
  name: gig-service
  namespace: production
spec:
  replicas: 3
  selector:
    matchLabels:
      app: gig-service
  template:
    spec:
      containers:
      - name: gig
        image: 19010853/ithust-gig:stable-1
        ports:
        - containerPort: 4003
        envFrom:
        - secretRef:
            name: backend-secrets
        resources:
          requests:
            memory: "256Mi"
            cpu: "250m"
          limits:
            memory: "512Mi"
            cpu: "500m"
\end{lstlisting}

\textbf{Service.yaml} --- Internal networking:
\begin{lstlisting}[style=BashStyle, caption={K8s Service}]
apiVersion: v1
kind: Service
metadata:
  name: gig-service
  namespace: production
spec:
  selector:
    app: gig-service
  ports:
  - port: 4003
    targetPort: 4003
  type: ClusterIP
\end{lstlisting}

\textbf{Ingress.yaml} --- External access qua Traefik:
\begin{lstlisting}[style=BashStyle, caption={K8s Ingress}]
apiVersion: networking.k8s.io/v1
kind: Ingress
metadata:
  name: ithust-ingress
  annotations:
    traefik.ingress.kubernetes.io/router.entrypoints: web,websecure
    traefik.ingress.kubernetes.io/router.tls.certresolver: letsencrypt
spec:
  rules:
  - host: api.ithustapp.com
    http:
      paths:
      - path: /api/v1/gigs
        backend:
          service:
            name: gig-service
            port: 4003
\end{lstlisting}

\subsection{Monitoring và Logging}

\textbf{Elastic APM} --- Application Performance Monitoring:
\begin{itemize}
    \item Theo dõi response times, error rates, throughput
    \item Distributed tracing qua các microservices
    \item Integration với Kibana để visualization
\end{itemize}

\textbf{Logging với Pino:}
\begin{lstlisting}[style=TypeScriptStyle, caption={Structured Logging}]
import { logger } from '@19010853/ithust-shared';

// Log levels: trace, debug, info, warn, error, fatal
logger.info('Order created successfully', {
  orderId: order._id,
  buyerId: order.buyerId,
  amount: order.price
});

// Logs được định dạng JSON, dễ parse bởi ELK stack
\end{lstlisting}

\section{Kết quả đạt được}

\subsection{Thống kê mã nguồn}

\begin{table}[H]
\centering
\begin{tabular}{|p{5cm}|p{9cm}|}
\hline
\textbf{Thông tin} & \textbf{Số liệu thống kê} \\ \hline
Số lượng microservices & 8 services (Node.js/TypeScript) \\ \hline
Frontend project & 1 project React/Vite/TypeScript \\ \hline
K8s manifests & 24+ YAML files cho deployments, services, ingress \\ \hline
Total lines of code & ~25,000+ lines (TypeScript/JavaScript) \\ \hline
Docker images & 9 images (8 services + 1 frontend) \\ \hline
GitHub Actions workflows & 8 workflows cho CI/CD \\ \hline
\end{tabular}
\caption{Thống kê mã nguồn ứng dụng}
\label{tab:thongke}
\end{table}

\subsection{Chức năng hoàn thành}

\begin{itemize}
    \item[\checkmark] \textbf{Xác thực}: Đăng ký, đăng nhập, JWT, quên mật khẩu
    \item[\checkmark] \textbf{User Management}: Profile buyer/seller, onboarding
    \item[\checkmark] \textbf{Gig Management}: CRUD gigs, media uploads
    \item[\checkmark] \textbf{Search}: Full-text search với Elasticsearch, filters
    \item[\checkmark] \textbf{Orders}: Shopping cart, checkout, order lifecycle
    \item[\checkmark] \textbf{Payments}: Stripe integration, webhooks
    \item[\checkmark] \textbf{Chat}: Real-time messaging với Socket.io
    \item[\checkmark] \textbf{Reviews}: Rating, feedback, seller response
    \item[\checkmark] \textbf{Notifications}: Email notifications qua RabbitMQ
    \item[\checkmark] \textbf{DevOps}: Docker, K3s, GitHub Actions CI/CD
\end{itemize}

\subsection{Minh họa chức năng (Ảnh cần chụp)}

\begin{figure}[H]
    \centering
    \fbox{\parbox{0.9\textwidth}{
        \centering
        \textbf{[ẢNH 1]}\\[0.2cm]
        \textbf{Trang chủ ITHUST App}\\[0.1cm]
        Giao diện tổng quan với categories, search bar, featured gigs
    }}
    \caption{Giao diện trang chủ}
    \label{fig:screenshot_home}
\end{figure}

\begin{figure}[H]
    \centering
    \fbox{\parbox{0.9\textwidth}{
        \centering
        \textbf{[ẢNH 2]}\\[0.2cm]
        \textbf{Kết quả tìm kiếm với Elasticsearch}\\[0.1cm]
        Hiển thị grid gigs, filters bên trái, sort options
    }}
    \caption{Giao diện tìm kiếm}
    \label{fig:screenshot_search}
\end{figure}

\begin{figure}[H]
    \centering
    \fbox{\parbox{0.9\textwidth}{
        \centering
        \textbf{[ẢNH 3]}\\[0.2cm]
        \textbf{Chi tiết Gig}\\[0.1cm]
        Gallery ảnh, mô tả chi tiết, packages/pricing, button ``Order Now''
    }}
    \caption{Giao diện chi tiết gig}
    \label{fig:screenshot_gig_detail}
\end{figure}

\begin{figure}[H]
    \centering
    \fbox{\parbox{0.9\textwidth}{
        \centering
        \textbf{[ẢNH 4]}\\[0.2cm]
        \textbf{Giỏ hàng và Thanh toán (Stripe)}\\[0.1cmcm]
        Order summary, Stripe payment form, card input
    }}
    \caption{Giao diện thanh toán}
    \label{fig:screenshot_checkout}
\end{figure}

\begin{figure}[H]
    \centering
    \fbox{\parbox{0.9\textwidth}{
        \centering
        \textbf{[ẢNH 5]}\\[0.2cm]
        \textbf{Chat Real-time}\\[0.1cm]
        Danh sách conversations, chat window, tin nhắn real-time
    }}
    \caption{Giao diện chat}
    \label{fig:screenshot_chat}
\end{figure}

\begin{figure}[H]
    \centering
    \fbox{\parbox{0.9\textwidth}{
        \centering
        \textbf{[ẢNH 6]}\\[0.2cm]
        \textbf{Dashboard Seller}\\[0.1cm]
        Quản lý gigs, orders, earnings, analytics
    }}
    \caption{Dashboard Seller}
    \label{fig:screenshot_seller_dashboard}
\end{figure}

\begin{figure}[H]
    \centering
    \fbox{\parbox{0.9\textwidth}{
        \centering
        \textbf{[ẢNH 7]}\\[0.2cm]
        \textbf{Kubernetes Dashboard / Kubectl}\\[0.1cm]
        Pods running, services, deployments trên K3s
    }}
    \caption{K3s Kubernetes deployment}
    \label{fig:screenshot_k8s}
\end{figure}

\begin{figure}[H]
    \centering
    \fbox{\parbox{0.9\textwidth}{
        \centering
        \textbf{[ẢNH 8]}\\[0.2cm]
        \textbf{RabbitMQ Management UI}\\[0.1cm]
        Queues, exchanges, message flow
    }}
    \caption{RabbitMQ Management Interface}
    \label{fig:screenshot_rabbitmq}
\end{figure}

\begin{figure}[H]
    \centering
    \fbox{\parbox{0.9\textwidth}{
        \centering
        \textbf{[ẢNH 9]}\\[0.2cm]
        \textbf{Kibana / Elastic APM}\\[0.1cm]
        Monitoring, logs, performance metrics
    }}
    \caption{Elastic Stack Monitoring}
    \label{fig:screenshot_kibana}
\end{figure}

% ============================================================
\chapter{KẾT LUẬN VÀ HƯỚNG PHÁT TRIỂN}
% ============================================================

\section{Kết luận}

Đồ án đã thành công trong việc nghiên cứu và xây dựng nền tảng \textbf{ITHUST App} --- một hệ thống microservices hoàn chỉnh cho giao dịch và trao đổi dịch vụ/công việc trực tuyến.

\textbf{Những điểm đạt được:}

\begin{itemize}
    \item \textbf{Kiến trúc hiện đại}: Hệ thống 8 microservices độc lập, mỗi service có trách nhiệm rõ ràng, giao tiếp qua REST API và message queue
    \item \textbf{Cloud-native}: Containerization với Docker, orchestration với Kubernetes (K3s), CI/CD tự động qua GitHub Actions
    \item \textbf{Tích hợp công nghệ đa dạng}: Node.js, React, MongoDB, PostgreSQL, MySQL, Redis, RabbitMQ, Elasticsearch, Socket.io, Stripe
    \item \textbf{High availability}: Horizontal scaling, health checks, auto-restart, fault tolerance qua circuit breakers và message queues
    \item \textbf{Real-time features}: Chat system với Socket.io, notification system với RabbitMQ
    \item \textbf{Search powerful}: Elasticsearch cho full-text search với relevance scoring
    \item \textbf{Security}: JWT authentication, bcrypt hashing, HTTPS, CORS, rate limiting
\end{itemize}

\textbf{Bài học kinh nghiệm:}
\begin{itemize}
    \item Kiến trúc microservice phù hợp cho hệ thống phức tạp, nhưng đòi hỏi chi phí operational cao hơn
    \item Event-driven architecture giúp giảm coupling nhưng tăng complexity trong việc đảm bảo consistency
    \item Kubernetes/K3s cung cấp nền tảng vững chắc cho deployment, nhưng cần thời gian học tập
    \item Observability (logging, monitoring, tracing) rất quan trọng trong distributed systems
\end{itemize}

\section{Hướng phát triển}

\subsection{Tính năng nghiệp vụ}

\begin{itemize}
    \item \textbf{Dispute Resolution}: Hệ thống xử lý tranh chấp giữa buyer và seller
    \item \textbf{Subscription/Membership}: Các gói dịch vụ premium cho sellers
    \item \textbf{Mobile App}: React Native app cho iOS/Android
    \item \textbf{AI-powered}: Recommendation engine, auto-tagging, sentiment analysis cho reviews
    \item \textbf{Multi-language}: i18n support cho thị trường quốc tế
\end{itemize}

\subsection{Công nghệ và kiến trúc}

\begin{itemize}
    \item \textbf{Service Mesh}: Istio hoặc Linkerd cho advanced traffic management, mTLS
    \item \textbf{GraphQL Federation}: Thay thế REST với GraphQL cho flexible queries
    \item \textbf{Event Sourcing}: Lưu trữ events thay vì state cho audit và replay
    \item \textbf{Saga Pattern}: Quản lý distributed transactions
    \item \textbf{Chaos Engineering}: Testing fault tolerance với tools như Chaos Monkey
    \item \textbf{GitOps}: ArgoCD cho declarative continuous delivery
\end{itemize}

\subsection{DevOps và Monitoring}

\begin{itemize}
    \item \textbf{Prometheus + Grafana}: Metrics collection và visualization
    \item \textbf{Jaeger}: Distributed tracing
    \item \textbf{ELK Stack hoàn chỉnh}: Centralized logging
    \item \textbf{Automated Testing}: Unit tests, integration tests, contract tests, E2E tests
    \item \textbf{Multi-environment}: Staging environment trước production
\end{itemize}

\vspace{1cm}
\begin{flushright}
    \textit{Xin chân thành cảm ơn!}
\end{flushright}

% ---- TÀI LIỆU THAM KHẢO ----
\addcontentsline{toc}{chapter}{TÀI LIỆU THAM KHẢO}
\begin{thebibliography}{99}

\bibitem{microservices_patterns}
Chris Richardson.
\textit{Microservices Patterns: With examples in Java}.
Manning Publications, 2018.

\bibitem{building_microservices}
Sam Newman.
\textit{Building Microservices: Designing Fine-Grained Systems}.
O'Reilly Media, 2021 (2nd Edition).

\bibitem{nodejs_design_patterns}
Mario Casciaro, Luciano Mammino.
\textit{Node.js Design Patterns}.
Packt Publishing, 2020 (3rd Edition).

\bibitem{kubernetes_up_running}
Kelsey Hightower, Brendan Burns, Joe Beda.
\textit{Kubernetes: Up and Running}.
O'Reilly Media, 2022 (3rd Edition).

\bibitem{rabbitmq_tutorials}
RabbitMQ Documentation.
\textit{RabbitMQ Tutorials}.
\url{https://www.rabbitmq.com/getstarted.html}, last visited May 2026.

\bibitem{elasticsearch_guide}
Elastic Documentation.
\textit{Elasticsearch Guide}.
\url{https://www.elastic.co/guide/en/elasticsearch/reference/current/index.html}, last visited May 2026.

\bibitem{socketio_docs}
Socket.io Documentation.
\textit{Socket.io Docs}.
\url{https://socket.io/docs/v4/}, last visited May 2026.

\bibitem{stripe_api}
Stripe Documentation.
\textit{Stripe API Reference}.
\url{https://stripe.com/docs/api}, last visited May 2026.

\bibitem{react_docs}
Meta Open Source.
\textit{React Documentation}.
\url{https://react.dev/}, last visited May 2026.

\bibitem{k3s_docs}
Rancher Labs.
\textit{K3s Documentation}.
\url{https://docs.k3s.io/}, last visited May 2026.

\bibitem{clean_architecture}
Robert C. Martin.
\textit{Clean Architecture: A Craftsman's Guide to Software Structure and Design}.
Prentice Hall, 2017.

\end{thebibliography}

\end{document}
